\begin{homeworkProblem}[1-23]
  Verifique quais dos sistemas lineares $x^{\prime}=Ax$, possuem uma integral primeira, no caso em que $A\in M(2)$ está na forma de Jordan. 
  \begin{solution}
  \begin{itemize}
    \item Para o caso
      \begin{equation}\label{eq:1-23-caso1}
        \begin{pmatrix}
          x^{\prime} \\
          y^{\prime}
        \end{pmatrix}=\begin{pmatrix}
          \lambda_1 & 0 \\
          0 & \lambda_2
        \end{pmatrix}\begin{pmatrix}
          x \\
          y
        \end{pmatrix}.
      \end{equation}
      Considere o seguinte método.\\
      Note que pela regra da cadeia temos que
      \begin{align*}
        \frac{dy}{dx}=\frac{y^{\prime}}{x^{\prime}}=\frac{\lambda_2y}{\lambda_1x}.
      \end{align*}
      Equação com solução
      \begin{align*}
        \lambda_1\ln(|y|)-\lambda_2\ln(|x|)=c.
      \end{align*}
      Assim, defina
      \begin{align*}
        H(x,y)=\lambda_1\ln(|y|)-\lambda_2\ln(|x|).
      \end{align*}
      Logo temos que se $(x,y)$ é solução de \eqref{eq:1-23-caso1}, então
      \begin{align*}
        H(x,y)=\lambda_1\ln(|y|)-\lambda_2\ln(|x|)=0.
      \end{align*}
      Isto é que $H$ é integral primeira de $A$.
    \item Para o caso
      \begin{equation}\label{eq:1-23-caso2}
        \begin{pmatrix}
          x^{\prime} \\
          y^{\prime}
        \end{pmatrix}=\begin{pmatrix}
          \lambda & 0 \\
          1 & \lambda
        \end{pmatrix}\begin{pmatrix}
          x \\
          y
        \end{pmatrix}.
      \end{equation}
      Considere o seguinte método.\\
      Note que pela regra da cadeia temos que
      \begin{align*}
        \frac{dy}{dx}=\frac{y^{\prime}}{x^{\prime}}=\frac{x+\lambda y}{\lambda x}.
      \end{align*}
      Logo, usando o fator integrante $e^{\int-\frac{1}{x}\,dx}=\frac{1}{x}$, podemos garantir que a sua solução é
      \begin{align*}
        \frac{y}{x}-\frac{\ln(|x|)}{\lambda}=c.
      \end{align*}
      Assim, defina
      \begin{align*}
        H(x,y)=\frac{y}{x}-\frac{\ln(|x|)}{\lambda}.
      \end{align*}
      Logo temos que se $(x,y)$ é solução de \eqref{eq:1-23-caso2}, então
      \begin{align*}
        H(x,y)=\frac{y}{x}-\frac{\ln(|x|)}{\lambda}=0.
      \end{align*}
      Isto é que $H$ é integral primeira de $A$.
    \item Para o caso
      \begin{equation}\label{eq:1-23-caso3}
        \begin{pmatrix}
          x^{\prime} \\
          y^{\prime}
        \end{pmatrix}=\begin{pmatrix}
          \Re \lambda & \Im \lambda \\
          -\Im \lambda & \Re \lambda
        \end{pmatrix}\begin{pmatrix}
          x \\
          y
        \end{pmatrix}.
      \end{equation}
      Considere o seguinte método.\\
      Note que levando a coordenadas polares, temos que
      \begin{align*}
        r^{\prime}=\Re\lambda r,\quad \theta^{\prime}=-\Im\lambda.
      \end{align*}
      Note que pela regra da cadeia temos que
      \begin{align*}
        \frac{dr}{d\theta}=\frac{r^{\prime}}{\theta^{\prime}}=-\frac{\Re\lambda r}{\Im\lambda}.
      \end{align*}
      Equação com solução
      \begin{align*}
        \ln(|r|)+\frac{\Re\lambda}{\Im\lambda}\theta=c.
      \end{align*}
      Assim, defina
      \begin{align*}
        H(x,y)=\frac{\ln(x^{2}+y^{2})}{2}+\frac{\Re\lambda}{\Im\lambda}\arctan\left( \frac{y}{x} \right).
      \end{align*}
      Logo temos que se $(x,y)$ é solução de \eqref{eq:1-23-caso3}, então
      \begin{align*}
        H(x,y)=\frac{\ln(x^{2}+y^{2})}{2}+\frac{\Re\lambda}{\Im\lambda}\arctan\left( \frac{y}{x} \right)=0.
      \end{align*}
      Isto é que $H$ é integral primeira de $A$.
  \end{itemize}
  \end{solution}
\end{homeworkProblem}
